

Many-query computations for parametrized incompressible flows arise in
aerodynamic design, flow control, optimization, uncertainty quantification,
inverse problems, and real-time prediction. In these settings, variations in
geometry, boundary conditions, material properties, forcing, or operating
parameters can affect not only the magnitude of the flow response but also
its qualitative temporal behavior. A parametrized flow may exhibit steady,
weakly oscillatory, or sustained periodic dynamics over different portions
of the parameter domain. High-fidelity numerical simulations can resolve
these behaviors accurately, but their computational cost becomes prohibitive
when solutions are required repeatedly for many parameter values or over long
time horizons. Reduced-order models (ROMs) provide an efficient alternative
by approximating the high-dimensional flow dynamics in a low-dimensional
representation. For broad-range parametrized problems, an effective ROM must
therefore provide not only accurate predictions within a given operating
regime, but also robust predictive performance as the underlying flow
dynamics change across the parameter domain.

Reduced-order models (ROMs) reduce this computational burden by representing the high-dimensional flow state in a low-dimensional subspace\cite{benner2015,hesthaven2016,quarteroni2016}. Broadly, ROMs can be divided into projection-based methods, data-driven approaches, and Physics-data combined. Among projection-based methods, proper orthogonal decomposition (POD) constructs an energy-optimal linear basis from flow snapshots \cite{sirovich1987,berkooz1993}, and Galerkin projection transfers the governing equations onto this basis. The resulting POD–Galerkin models have been widely used for transient and post-transient flows \cite{noack2003,rowley2004,roz2007}. However, their predictive accuracy can degrade when truncation removes modes that contribute significantly to dissipation, nonlinear interactions, or pressure coupling. In such cases, the reduced dynamics may exhibit phase drift, biased mean states, and inaccurate pressure fields, even when the POD projection error itself is small. Stabilization, calibration, and closure modeling can alleviate these
deficiencies \cite{amsallem2012}, but their effectiveness remains closely
linked to the reduced representation and dynamical regime for which they are
constructed.

Data-driven ROMs provide a complementary approach by learning reduced
dynamics directly from simulation or experimental trajectories. Dynamic mode
decomposition and Koopman-based methods identify evolution in observable
coordinates \cite{schmid2010,tu2014,williams2015}, sparse identification
constructs parsimonious nonlinear models
\cite{brunton2016,champion2019}, and operator inference estimates reduced
operators non-intrusively from trajectory data \cite{peherstorfer2016}.
Neural approaches further extend this class through coefficient-space
regression, recurrent architectures, and neural operators
\cite{hesthaven2018,pathak2018,vlachas2018,lu2021,li2021,kovachki2023}.
These approaches increase the flexibility with which reduced dynamics can be
approximated, although accurate local or one-step prediction does not by
itself guarantee stable autonomous evolution. Errors introduced during
prediction are recursively propagated and may gradually move the reduced
state away from the trajectories represented in the training data
\cite{duraisamy2019,brunton2020}.These complementary strengths and limitations motivate hybrid formulations that combine physical projection with data-driven modeling.

Physics data combined reduced models seek to combine the structural information retained by
projection-based methods with the approximation capability of learned
corrections. Neural closures have been introduced to compensate for
unresolved projection errors \cite{dar2023}, while nonlinear and learned
manifolds provide more expressive reduced representations when linear
subspaces become inefficient \cite{lee2020,fresca2022,geelen2023}.
Domain-decomposed nonlinear manifolds, structure-informed latent models, and
local non-intrusive approximations provide related strategies for complex
parametrized systems \cite{diaz2024,park2024,discacciati2024}. These
developments demonstrate that physical projection and data-driven modeling
can be combined effectively: projected equations retain known physical
structure, while learned components compensate for unresolved reduced
dynamics. Nevertheless, improving the predictive accuracy of a reduced model
within a prescribed representation does not by itself resolve the additional
difficulty introduced when the qualitative flow regime changes across the
parameter domain.

The central issue is therefore not merely predictive accuracy, but representation itself. The reduced basis, scaling, and dynamical closure appropriate for one regime may be inappropriate for another.  Let $\mathcal{M}$ denote the set of flow states generated
over all admissible parameters and times. Although the full discrete state
belongs to a linear space, $\mathcal{M}$ is generally nonlinear, and its
local geometry may vary substantially with the parameter. For a system
passing through steady, Hopf-onset, and developed periodic dynamics, the
corresponding state subsets can exhibit different dominant spatial
directions and markedly different temporal characteristics. Denoting these
subsets by $\mathcal{M}_S$, $\mathcal{M}_H$, and $\mathcal{M}_P$, one has
\[
  \mathcal{M}
  =
  \mathcal{M}_S
  \cup
  \mathcal{M}_H
  \cup
  \mathcal{M}_P .
\]
The Kolmogorov width of this union can decay more slowly than those of the
individual regime-restricted subsets, so that a single linear approximation
space may require substantially more modes to represent the complete
parameter range. Moreover, the associated reduced dynamics must accommodate
qualitatively different temporal behaviors, including relaxation toward a
steady state, oscillatory growth near onset, and nonlinear saturation in the
periodic regime. Hence, robustness across regime changes involves both the
reduced representation and the reduced dynamics, rather than the accuracy of
either component in isolation.

Localization offers a natural means of addressing this regime dependence.
Local reduced spaces and cluster-based models approximate different portions
of the solution set with specialized low-dimensional representations
\cite{kaiser2014,li2021cluster}, while mixture-of-experts architectures
provide a general mechanism for assigning specialized predictors to
different regions of an input space
\cite{jacobs1991,jordan1994,shazeer2017,fedus2022}. For parametrized ROMs,
however, localization introduces an additional mathematical challenge.
Local models constructed from different portions of the solution manifold
generally possess different means, POD bases, coefficient scalings, and
reduced operators. Their reduced coordinates therefore do not admit a
canonical componentwise correspondence. A localized ROM must consequently
address two coupled questions: how to select an appropriate regime-local
model as the parameter varies, and how to combine predictions from
neighboring local models without assuming that their reduced coordinates
belong to a common latent space.



```latex
Motivated by these considerations, we propose the
\emph{Regime-Adaptive Localized Mixture-of-Experts Reduced-Order Model}
(RAL-MoE-ROM), a hierarchical multi-chart framework for parametrized flows
with regime-dependent dynamics. Rather than describing the complete
parameter domain through a single reduced representation, RAL-MoE-ROM
covers the solution manifold by overlapping regime-local POD charts. Each
chart defines its own reduced coordinates and projected operators and is
equipped with an independently evolved physics--data specialist.

Within each specialist, the projected reduced dynamics provide a
physics-based backbone, while the unresolved contribution is represented by
a structured sparse mixture of experts. The internal MoE employs hierarchical
routing to activate only a subset of specialized expert maps for a given
reduced state and its recent history. Each expert combines nonlinear,
affine, and low-rank bilinear components, providing a structured
approximation of unresolved interactions while retaining the information
contained in the projected reduced operator. This construction introduces
adaptivity within each dynamical regime without replacing the underlying
reduced model by an unconstrained black-box predictor.

A second level of adaptivity operates between regime-local specialists.
Because independently constructed local charts generally possess
noncommensurate reduced coordinates, RAL-MoE-ROM does not combine their
modal states or reduced operators directly. Instead, an outer hierarchical
mechanism identifies the relevant neighboring specialists according to the
regime structure and their dynamical admissibility. In transition regions,
the admissible local models are advanced independently in their native
reduced coordinates, and the proposed Top-2 convex correction (T2-C)
combines their reconstructed physical states in the common physical space.
The resulting architecture therefore separates intra-regime dynamical
specialization from inter-regime model selection and coupling.

The principal methodological contributions are summarized as follows:
\begin{itemize}

\item \textbf{A regime-adaptive multi-chart reduced representation.}
The parametrized solution manifold is represented by overlapping local POD
charts associated with qualitatively different dynamical regimes. Each chart
retains an independent reduced coordinate system and projected dynamics,
allowing the reduced representation to adapt to changes in the local
geometry and temporal behavior of the solution set.

\item \textbf{A structured sparse mixture-of-experts formulation for
regime-local reduced dynamics.}
Each local chart is augmented by an autonomous physics--data specialist in
which a hierarchical sparse router selects among structured expert
corrections. The expert maps combine nonlinear, affine, and low-rank
bilinear components, so that unresolved reduced dynamics are modeled through
a conditional architecture that retains both general nonlinear flexibility
and explicit low-order interaction structure. This establishes a second
level of specialization within each regime-local ROM.

\item \textbf{An admissibility-constrained assembly of noncommensurate
local ROMs.}
The regime-local specialists are coupled without introducing a common latent
coordinate system. The outer hierarchy restricts interaction to neighboring
regimes and admits soft combination only when the corresponding local models
are dynamically applicable. The T2-C mechanism then combines independently
evolved local predictions after reconstruction in the common physical space,
thereby avoiding direct interpolation of modal states or reduced operators.

\end{itemize}
```


The remainder of the paper is organized as follows.
Section~\ref{sec:regime_local_roms} introduces the regime-local reduced
representations and projected operators.
Section~\ref{sec:regime_local_specialists} develops the physics--data
specialists.
Section~\ref{sec:hierarchical_ral_moe} presents the hierarchical
cross-chart assembly and T2-C fusion.
The subsequent section reports the numerical experiments, and the final
section summarizes the main findings and limitations.
```





















